\documentclass{beamer}
%% fink install texlive
\usetheme{Warsaw}
\usepackage{amsmath,amssymb}
\usepackage{stmaryrd}
\usepackage{graphicx}
%\usepackage{haskell}
\usepackage{code}
%\usepackage{proof}
\usepackage{theorem}
\usepackage{pstricks} 
\usepackage{listings}

\include{macros-ms}

\newcommand{\rev}[1]{{#1}^{\mbox{\scriptsize r}}}
\newcommand{\lquo}{\backslash}
\newcommand{\rquo}{/}
\newcommand{\diff}{-}
\newcommand{\isect}{\cap}
\newcommand{\mymid}{~\redtxt{\mid}~}
\newcommand{\lconf}[2]{#1:{\tt #2}}

\newcommand{\ignore}[1]{}

\newcommand{\magtxt}[1]{{\magenta #1}}
\newcommand{\redtxt}[1]{{\red #1}}
\newcommand{\bluetxt}[1]{{\blue #1}}
\newcommand{\greytxt}[1]{{\gray #1}}
\newcommand{\greentxt}[1]{{\green #1}}
\newcommand{\mmleq}{\leq}
\newcommand{\pow}{\^{}}
\newcommand{\venv}{\Delta}
\newcommand{\mleq}{\mbox{\tt leq}}
\newcommand{\mas}{\mbox{\tt as}}
\newcommand{\subt}{<\!:}

\newcommand{\Nturns}{\, \vdash_{\mbox{\tiny lnf}} \,}

\newcommand{\deriv}[2]{#1\backslash #2}
\newcommand{\pderiv}[2]{pd_{#2}(#1)} %% {#1\backslash_p #2}
\newcommand{\pderivFail}[2]{pd^{err}_{#2}(#1)} %% {#1\backslash_p #2}
\newcommand{\mkEmp}[1]{{#1}_{\epsilon}}
\newcommand{\strip}[1]{#1\downarrow}
\newcommand{\beb}{\begin{exampleblock}}
\newcommand{\eeb}{\end{exampleblock}}
\newenvironment{ttprog}{\begin{trivlist}\item \tt
        \begin{tabbing}}{\end{tabbing}\end{trivlist}}


\newenvironment{grammar}{%
         \begin{center} \small%
         $\begin{array}{rcll}
         }{%
         \end{array}$\end{center}\ignorespaces%
         }

\newcommand{\rem}[1]{}
\newcommand{\comment}[1]{}

\newcommand{\kl}[1]{}

\bibliographystyle{plainnat}

\begin{document}
%\lstset{language=python}
\logo{
\includegraphics[height=0.8cm]{pics/nyp-logo.eps}\vspace{230pt}
\includegraphics[height=0.6cm]{pics/sit-logo.png}\vspace{230pt}
}
\title{Scala Parallel and Concurrent Programming} 
\author{
 Kenny Zhuo Ming Lu
}
\institute[Inst.]{School of Information Technology, Nanyang Polytechnic}
\date{\today} 


\frame{\titlepage} 
%-------------------------------------------------------------------
%-------------------------------------------------------------------
\begin{frame}[fragile]
\frametitle{Data Parallelism}
Scala has built-in support of data parallelism. For instance,
\beb{}
\begin{code}
object Fib extends App {
  def fib(x:Int):Int = x match {
    case 0 => 0
    case 1 => 1
    case n => fib(n-1) + fib(n-2)
  }
  println((30 to 40).toList.map(fib))
}
\end{code}
\eeb
\begin{code}
$ time scala Fib
List(832040, 1346269, 2178309, ... )
real	0m1.621s
user	0m1.624s
sys	0m0.059s
\end{code}
%$
\end{frame}

%-------------------------------------------------------------------
%-------------------------------------------------------------------
\begin{frame}[fragile]
\frametitle{Data Parallelism}
{\tt .par} converts a sequentials data collection into a parallel
counter part. {\tt .seq} has the opposite effect.
\beb{}
\begin{code}
object ParFib extends App {
  def fib(x:Int):Int = x match {
    case 0 => 0
    case 1 => 1
    case n => fib(n-1) + fib(n-2)
  }
  println((30 to 40).toList.par.map(fib))
}
\end{code}
\eeb
\begin{code}
$ time scala ParFib
List(832040, 1346269, 2178309, ... )
real	0m0.891s
user	0m1.901s
sys	0m0.061s
\end{code}
%$
\end{frame}


%-------------------------------------------------------------------
%-------------------------------------------------------------------
\begin{frame}[fragile]
\frametitle{Concurrent Programming}

\begin{itemize}
\item Multi-threading
\item Actors
\item Async computation via Future
\end{itemize}

\end{frame}

%-------------------------------------------------------------------
%-------------------------------------------------------------------
\begin{frame}[fragile]
\frametitle{Concurrent Programming}

Multi-threading in Scala is similar to Java
\beb{}
\begin{code}
  object ThreadFib extends App {
    def fib(x:Int):Int = x match {
      case 0 => 0
      case 1 => 1
      case n => fib(n-1) + fib(n-2)
    }

    for (i <- 30 to 40) { 
      new Thread(new Runnable {
        def run() {
          println(fib(i))
        }
      }).start
    }
  }
\end{code}
\eeb
\end{frame}


%-------------------------------------------------------------------
%-------------------------------------------------------------------
\begin{frame}[fragile]
\frametitle{Actors}
Actor model is a concurrent programming model. 
\begin{itemize}
 \item Actors are primitive objects with specific roles and functionalities
 \item Actors (objects) are coordinating with each other via message passing
\end{itemize}
\end{frame}

%-------------------------------------------------------------------
%-------------------------------------------------------------------
\begin{frame}[fragile]
\frametitle{Actors}
Let's consider a ping pong simulation using Scala actors

\begin{itemize}
\item There are two actors {\tt Ping} and {\tt Pong}. 
\item {\tt Ping} keeps an internal counter
   \begin{itemize}
      \item when it receives a {\tt start} message, 
        it increments the counter and send a {\tt ping} message to {\tt Pong}
       \item when it receives a {\tt pong} message, it increments the
         counter. 
           \begin{itemize} 
               \item if the counter $>$ 99, it sends a {\tt stop}
                 message to {\tt pong} then terminates.
               \item  otherwise, sends a {\tt ping} message to {\tt Pong}.
           \end{itemize}
   \end{itemize}
\item {\tt Pong} 
  \begin{itemize}
    \item when it receives a {\tt stop} message, it terminates.
    \item when it receives a {\tt ping} message, it sends a {\tt pong}
      message to {\tt Ping}.
  \end{itemize}
\end{itemize}

\beb{}
\begin{code}
case object PingMessage
case object PongMessage
case object StartMessage
case object StopMessage
\end{code}
\eeb
\end{frame}


%-------------------------------------------------------------------
%-------------------------------------------------------------------
\begin{frame}[fragile]
\frametitle{Actors}
\beb{}
\begin{code}
class Ping(pong: ActorRef) extends Actor {
  var count = 0
  def incrementAndPrint = {count += 1; println("ping")}
  def receive = {
    case StartMessage => 
      incrementAndPrint
      pong ! PingMessage
    case PongMessage => 
      incrementAndPrint
      if (count > 99) {
        sender ! StopMessage
        println("ping stopped")
        context.stop(self)
      } 
      else { sender ! PingMessage }
  }
}
\end{code}
\eeb
\end{frame}

%-------------------------------------------------------------------
%-------------------------------------------------------------------
\begin{frame}[fragile]
\frametitle{Actors}

\beb{}
\begin{code}
class Pong extends Actor {
  def receive = {
    case PingMessage =>
      println("  pong")
      sender ! PongMessage
    case StopMessage =>
      println("pong stopped")
      context.stop(self)
  }
}
\end{code}
\eeb
\end{frame}
%-------------------------------------------------------------------
%-------------------------------------------------------------------
\begin{frame}[fragile]
\frametitle{Actors}

\beb{}
\begin{code}

object PingPongTest extends App {
  val system = ActorSystem("PingPongSystem")
  val pong = system.actorOf(Props[Pong], 
  name = "pong")
  val ping = system.actorOf(Props(new Ping(pong)), 
  name = "ping")
  ping ! StartMessage
}
\end{code}
\eeb
\end{frame}

%-------------------------------------------------------------------
%-------------------------------------------------------------------
\begin{frame}[fragile]
\frametitle{Actors}
The ping pong example is pretty contrived, let's consider a less
contrived example in which we try to compute {\tt fib(n)} in
parallel. The idea is as follows,
\begin{itemize}
\item if n is less than 10, we compute fib(n) locally
\item if n is greater than or equal to 10, we create two fib
  actors. The first actor computes {\tt fib(n-1)} and the second actor
  computes {\tt fib(n-2)}. The two actors can run in parallel if there
  is enough threads. When both actors revert with some results, we sum up the results.
\end{itemize}
Let's define the messages for communication as follows,
\beb{}
\begin{code}
  case class FibRequest(n:Int)
  case class FibResult(n:Int)
  case object Stop
\end{code}
\eeb
\end{frame}

%-------------------------------------------------------------------
%-------------------------------------------------------------------
\begin{frame}[fragile]
\frametitle{Actors}
\beb{}
{\scriptsize
\begin{code}
class FibActor extends Actor {
  def receive = {
    case Stop => context.stop(self)
    case FibRequest(n) if n < 10 => sender ! FibResult(fib(n))
    case FibRequest(n) if n >= 10 => {
      val child1 = context.actorOf(Props[FibActor]) 
      val child2 = context.actorOf(Props[FibActor])
      val future1:Future[Any] = child1 ? FibRequest(n-1) // ``ask'' op
      val future2:Future[Any] = child2 ? FibRequest(n-2)
      val dest = sender
      val future1and2 = for ( // Future[T] is a monad
        FibResult(v1) <- future1;
        FibResult(v2) <- future2 
      ) yield v1+v2
      future1and2 onComplete {
        case Success(v) => 
        child1 ! Stop; child2 ! Stop
        dest ! FibResult(v) // can't use sender
        case Failure(e) => println("failed!" + e.toString)
      }
    }
  }
}
\end{code}
}
\eeb
\end{frame}

%-------------------------------------------------------------------
%-------------------------------------------------------------------
\begin{frame}[fragile]
\frametitle{Actors}
\beb{}
\begin{code}
def main(args:Array[String]):Unit = {
  val system = ActorSystem("FibSystem")
  val fibActor = system.actorOf(Props[FibActor])
  val n = if (args.length > 0) args(0).toInt else 10
  val f = fibActor ? FibRequest(n)
  f onComplete {
    case Success(v) => { println(v) ; fibActor ! Stop}
    case Failure(e) => println("failed!" + e.toString)
  }
}
\end{code}
\eeb
\end{frame}

\end{document}
